\begin{abstract}
A shared library for finite stochastic dynamics using normalized real weights,
finite sums, and limits of finite-time probabilities. These are generic library
results, rather than statements attributed to a particular opinion-dynamics paper.
\end{abstract}

\section{Uniform expectations}
\begin{definition}[Finite uniform average]\label{def:avg}
\lean{Dynamics.avg}\leanok
For a finite type $A$, set $\operatorname{avg}(f)=\sum_{a\in A}f(a)/|A|$.
When $A$ is empty this is zero, following Lean's division convention.
\end{definition}
\begin{lemma}[Product independence]\label{lem:uniform-product}
\lean{Dynamics.avg_prod_pi,Dynamics.avg_equiv}\leanok\uses{def:avg}
Uniform averages factor over independent coordinates and are invariant under
bijections of the sample space.
\end{lemma}
\begin{lemma}[Finite trajectories]\label{lem:uniform-trajectory}
\lean{Dynamics.expList,Dynamics.expList_eq_avg_ofFn}\leanok\uses{def:avg}
The recursively defined expectation over $T$ independent draws equals the
uniform average over functions $\{0,\ldots,T-1\}\to A$.
\end{lemma}
\begin{proof}\leanok
Separate the first coordinate using the finite-function/product equivalence,
then induct on $T$. The empty sample space is handled separately.
\end{proof}
\begin{lemma}[Averages commute]\label{lem:avg-commute}
\lean{Dynamics.avg_avg_swap,Dynamics.expList_avg_comm,Dynamics.expList_zero_fun,Dynamics.expList_finset_sum}\leanok\uses{lem:uniform-trajectory}
Iterated uniform averages over two finite sets commute; in particular an outer
average commutes with the trajectory expectation $\mathrm{expList}$, which also
commutes with finite sums and vanishes on the zero functional.
\end{lemma}
\begin{proof}\leanok
Exchange the order of finite summation; for $\mathrm{expList}$, induct on $T$.
\end{proof}

\section{Weighted distributions and kernels}
\begin{definition}[Normalized real weights]\label{def:distribution}
\lean{Dynamics.Distribution,Dynamics.Distribution.expect,Dynamics.Distribution.prob}\leanok
A distribution consists of nonnegative real weights $p_a$ with $\sum_a p_a=1$.
Its expectation is $\sum_a p_a f(a)$, and an event probability is the expectation
of its indicator. Such a distribution cannot exist on an empty type.
\end{definition}
\begin{lemma}[Uniform agreement]\label{lem:uniform-agreement}
\lean{Dynamics.Distribution.uniform_expect}\leanok\uses{def:distribution,def:avg}
On a nonempty finite type the normalized uniform distribution agrees with the
existing finite average.
\end{lemma}
\begin{lemma}[Independent products and pushforward]\label{lem:weighted-product}
\lean{Dynamics.Distribution.independent_expect_prod,Dynamics.Distribution.independent_expect_eval,Dynamics.Distribution.map_expect}\leanok\uses{def:distribution}
Independent coordinates have product weights. Expectations of product
observables factor, and transporting weights along a function commutes with expectation.
\end{lemma}
\begin{proof}\leanok
Distribute products over finite sums; for pushforward interchange the two sums
and evaluate each fiber indicator.
\end{proof}
\begin{definition}[Transition kernel]\label{def:kernel}
\lean{Dynamics.Kernel,Dynamics.Kernel.iterate,Dynamics.Kernel.trajectory}\leanok\uses{def:distribution}
A finite kernel assigns a distribution to every state. Its operator iterates
compute expectations at finite times; trajectory functionals additionally
record the history.
\end{definition}
\begin{lemma}[Endpoint agreement]\label{lem:endpoint}
\lean{Dynamics.Kernel.trajectory_endpoint}\leanok\uses{def:kernel}
Ignoring history in a trajectory expectation recovers the iterated transition operator.
\end{lemma}

\section{Stationarity and absorption}
\begin{theorem}[Existence of stationary weights]\label{thm:stationary}
\lean{Dynamics.Kernel.exists_stationary}\leanok\uses{def:kernel}
Every stochastic kernel on a nonempty finite type has a distribution $p$ with $pK=p$.
\end{theorem}
\begin{proof}\leanok
Average the first $n+1$ evolved distributions. The stationarity defect of this
Ces\`aro average is $(p_{n+1}-p_0)/(n+1)$, which tends to zero coordinatewise.
Compactness of the finite probability simplex gives a convergent subsequence;
continuity of finite sums implies that its limit is stationary.
\end{proof}
\begin{lemma}[Uniform absorption block]\label{lem:block}
\lean{Dynamics.Kernel.exists_uniform_block}\leanok\uses{def:kernel}
Let $f$ be a zero-one survival indicator with $Kf\le f$. Suppose for each
state $a$ there is $n_a$ with $(K^{n_a}f)(a)<1$. There are $m>0$ and
$0\le q<1$ such that $K^m f\le qf$.
\end{lemma}
\begin{proof}\leanok
Take one plus the maximum of the finitely many access times. Antitonicity of
survival probabilities puts every state below one at that time. Their finite
maximum supplies $q$. At states with $f=0$, absorption keeps the expectation zero.
\end{proof}
\begin{theorem}[Geometric decay and convergence]\label{thm:absorption}
\lean{Dynamics.Kernel.geometric_blocks,Dynamics.Kernel.finite_absorption}\leanok\uses{lem:block}
Under these hypotheses, $K^{nm}f\le q^n f$, and $K^t f(a)\to0$ for every state $a$.
\end{theorem}
\begin{proof}\leanok
Induct over blocks using positivity and linearity. The geometric bound tends
to zero. For arbitrary $t$, compare with $m\lfloor t/m\rfloor$ by antitonicity.
\end{proof}
